\scanpage{11}\mspage{\handwritten{10}}

qu'i\unsure{l|h} conviendra de poser les définitions adéquates). - Peut-être \unsure{Quillen} a-t-il une définition semi-si\unsure{mp|m}pliciale plus directe, sans passer par des $n$-catégories, des $K_i$ et plus précisément du H-espace qui devrait correspondre à $C$.\margin{\handwritten{\(\Vert\)\ \uuline{\unsure{; oui|si oui|oui}}}}
% slanted
\deleted{Dans le cas} Ce n'est certainement \inserted{en général} pas l'ensemble semi-simplicial $SC$ \deleted{engendré \unsure{à}} associé à $C$, puisque dans le cas d'une catégorie additive, $C$ est homotopiquement trivial (en effet le foncteur identique est ``\unsure{homotope}'' au foncteur nul, puisqu'il suffit d'avoir un homomorphisme d'un foncteur dans l'autre; d'où une équivalence d'homotopie entre $C$ et la catégorie nulle). Pourtant, il semble bien que dans le cas où $C$ est une \unsure{Gr-catégorie} i.e. admet des inverses, alors ce soit bien la construction ``naïve'' $SC$ qui est la bonne, puisqu'on trouve bien un $\handwritten{\pi}_{\unsure{o|0}}$ et un $\handwritten{\pi}_1$ qui sont aussi les $\handwritten{\pi}_o$ et $\handwritten{\pi}_1$ au sens catégorie sans Gr-structure.\margin{\handwritten{? \(\vert\)}}

\margin{\handwritten{\unsure{??|?} \(\vert\)}}\handwritten{À expliquer : pour des Gr-$n$-catégories, les $\pi_i$ au sens semi-simplicial sont aussi les $\pi_i$ au sens Gr-$n$-catégories, savoir les objets à isom près dans divers $\underline{\mathrm{Hom}}(1,1)$ ...... \deleted{Il y a aussi une interprétation de l'opération \unsure{du|des} $\unsure{\pi_1|\pi_0|\Omega\mathpunct{:}}$ \unsure{des|sur} \unsure{autres|lacets} en termes d'opération de shift géométrique évidente sur les $n$-catégories.....}}

% archive sheet number: 10
